\chapter{Transformer Mechanics}
\label{ch:transformer-mechanics}

\abstract{This chapter derives the \transformer from tensor shapes. It explains attention, masking, residual streams, normalization, feed-forward layers, positional information, and the implementation invariants needed to debug a model from first principles.}

\section{Chapter Spine}
The \transformer is a sequence-to-sequence computation graph whose core operation is content-dependent mixing. The reader should leave this chapter able to predict every major tensor shape in a forward pass and to explain why masking, normalization, and residual connections are necessary.

\section{Core Sections}
\subsection{Embedding and Positional Information}
Token embeddings, learned and sinusoidal positions, rotary embeddings, and where position enters attention.

\subsection{Scaled Dot-Product Attention}
Queries, keys, values, attention scores, softmax, dropout, and numerical stability.

\subsection{Masks and Causality}
Padding masks, causal masks, attention bias, and encoder-decoder cross-attention.

\subsection{Residual Streams and Feed-Forward Blocks}
Layer normalization, pre-norm versus post-norm, MLP expansion, gated activations, and residual information flow.

\section{Planned Exercises}
Implement single-head attention, multi-head attention, and a masked decoder block with shape assertions at every step.
